\chapter*{符号索引}
\markboth{符号索引}{符号索引}
\addcontentsline{toc}{chapter}{符号索引}
\sourcepage{247--248}
索引中的页码均为原书页码，可配合正文的原书页码提示查找。
\section*{一般记号}
\begin{longtable}{p{0.26\textwidth}p{0.66\textwidth}}
$\mathbb N$ & 自然数集（包括0）。\\
$\Z$ & 整数环。\\
$\R$ & 实数域。\\
$\R^*$ & $\{x\in\R:x\ne0\}$。\\
$\R_+$、$\R_-$ & 分别为 $\{x\in\R:x\ge0\}$、$\{x\in\R:x\le0\}$。\\
$\R_+^*$ & $\{x\in\R:x>0\}$。\\
$\C$、$\C^*$ & 复数域；$\{z\in\C:z\ne0\}$。\\
$\operatorname{Re}z$、$\operatorname{Im}z$ & 复数 $z$ 的实部、虚部。\\
$\PP^1$ & $\C\cup\{\infty\}$，黎曼球面，3。\\
$\PP^N$ & $N$ 维射影空间，142。\\
$[a,b]$ & $\{x\in\R:a\le x\le b\}$，其中 $a,b\in\R$，$a\le b$。\\
$]a,b]$ & $\{x\in\R:a<x\le b\}$，其中 $a,b\in\R$，$a<b$。\\
$M(m\times n,R)$ & 系数在 $R$ 中的 $m\times n$ 矩阵全体。\\
$\mathrm{GL}(n,R)$ & 系数在环 $R$ 中的可逆 $n\times n$ 矩阵群。\\
$A\Subset B$ & $A$ 为 $B$ 的相对紧子集。\\
$\partial A$ & $A$ 的边界。\\
\end{longtable}
\section*{层与函数空间}
\begin{longtable}{p{0.26\textwidth}p{0.66\textwidth}}
$\mathcal C$ & 连续函数层，41。\\
$\mathcal E$ & 可微函数层，60。\\
$\mathcal E^{(1)}$、$\mathcal E^{1,0}$、$\mathcal E^{0,1}$ & 63。\\
$\mathcal E^{(2)}$ & 65。\\
$\OO$ & 全纯函数层，41。\\
$\OO^*$、$\OO_D$、$\mathrm{GL}(n,\OO)$ & 分别见41、128、220。\\
$\Omega$、$\overline\Omega$、$\Omega_D$ & 分别见63、153、134。\\
$\MM$ & 亚纯函数层，41。\\
$\MM^*$、$\MM^{(1)}$ & 分别见41、65。\\
$\mathrm{Harm}^1$ & 154。\\
$|\mathcal F|$ & 与预层相伴的拓扑空间，42。\\
$\mathcal F_x$ & 预层在 $x$ 处的茎，41。\\
$\rho_x(f)$ & $f$ 在 $x$ 处的芽，42。\\
$\C\{z-a\}$ & 在 $a$ 附近收敛的泰勒级数，42。\\
$\C\{\{z-a\}\}$ & 在 $a$ 附近主部有限的洛朗级数域，57。\\
$\mathfrak m_a$、$\mathfrak m_a^2$ & 60。\\
$L^2(D,\OO)$ & 平方可积全纯函数空间，109。\\
$\mathcal D(X)$ & 紧支集可微函数空间，190。\\
$\mathcal D'(X)$ & 分布空间，190。\\
\end{longtable}
\section*{上同调理论}
\begin{longtable}{p{0.38\textwidth}p{0.54\textwidth}}
$B^1(\mathfrak U,\mathcal F)$、$B^1(G,A)$ & 分别见97、215。\\
$C^q(\mathfrak U,\mathcal F)$ & 96。\\
$C^q_{L^2}(\mathfrak U,\OO)$、$C_1(X)$ & 分别见112、161。\\
$H^1(\mathfrak U,\mathcal F)$、$H^1(X,\mathcal F)$ & 分别见97、99。\\
$H_1(X)$、$H^1(G,A)$ & 分别见161、215。\\
$Z^1(\mathfrak U,\mathcal F)$、$Z^1_{L^2}(\mathfrak U,\OO)$ & 分别见97、112。\\
$Z_1(X)$、$Z^1(G,A)$ & 分别见161、215。\\
$t^{\mathfrak U}_{\mathfrak V}$、$\mathfrak V\Subset\mathfrak U$ & 分别见98、112。\\
\end{longtable}
\section*{其他记号}
\begin{longtable}{p{0.38\textwidth}p{0.54\textwidth}}
$T_a^{(1)}$、$T_a^{1,0}$、$T_a^{0,1}$、$T_a^{(2)}$ & 分别见60、62、62、65。\\
$d$、$d'$、$d''$；$\partial$、$\overline\partial$ & 61、66；60。\\
$\pi_1(X,a)$、$\pi_1(X)$ & 基本群，17。\\
$\operatorname{cl}(u)$、$u^-$ & 分别见17、15。\\
$b_1(X)$ & 第一贝蒂数，157。\\
$\operatorname{Deck}(Y/X)$ & 覆盖变换群，33。\\
$\operatorname{Div}(X)$ & 除子群，127。\\
$\operatorname{Div}_p(X)$、$\operatorname{Div}_0(X)$ & 170。\\
$\deg(D)$、$(f)$、$(\omega)$ & 127。\\
$\operatorname{Jac}(X)$ & 雅可比簇，170。\\
$\wp$ & 魏尔斯特拉斯 $\wp$ 函数，13。\\
$\operatorname{Per}(\omega_1,\ldots,\omega_g)$ & 周期格，167。\\
$\operatorname{Pic}(X)$、$\operatorname{Pic}_0(X)$ & 皮卡群，170。\\
$\operatorname{Rh}^1(X)$、$\operatorname{Rh}_{\mathrm{hol}}^1(X)$ & 德拉姆群，125；全纯德拉姆群，206。\\
$\operatorname{Res}$ & 留数，63、132。\\
$\operatorname{Reg}(Y)$ & 179。\\
$\operatorname{Supp}$ & 函数或微分形式的支集，77、237。\\
\end{longtable}
